\section{Strict and stable finite palette bounds}\label{sec:finite}
We return to the triangular and nontriangular sectors of Section~\ref{sec:palettes} and assume $\delta>1/3$. The optimal matching of Section~\ref{sec:matching} has $M=2\nu$, and its marks satisfy the degree, root, and union hypotheses of Theorem~\ref{thm:polynomial}.

\begin{lemma}[Weighted triangle positivity]\label{lem:mantel}
If $m>1/4$, then $\Delta>0$.
\end{lemma}
\begin{proof}
If $\Delta=0$, then $Q$ is triangle-free, since all vertex weights are positive. Hence $D(u)+D(v)\le1$ on every edge, and summing with edge weights yields
\[
 \sum_vx_vD(v)^2
 =\sum_{uv\in E(Q)}x_ux_v\bigl(D(u)+D(v)\bigr)\le m.
\]
On the other hand, by the Cauchy--Schwarz inequality and $\sum_vx_vD(v)=2m$, the left side is at least $4m^2$. Thus $m\le1/4$, a contradiction.
\end{proof}

\begin{theorem}[Finite palette bound]\label{thm:finite}
Every positively weighted finite simple graph with $\delta>1/3$ and $m>1/4$ satisfies
\begin{equation}\label{eq:finite}
 \Phi\ge\frac32m-\frac14.
\end{equation}
\end{theorem}
\begin{proof}
By Lemma~\ref{lem:mantel}, $\Delta>0$. Theorem~\ref{thm:polynomial} and Lemma~\ref{lem:resources} give
\[
 \Phi_F\ge\max_wq(w)\ge\frac B\Delta
 \ge f+\nu+\frac m2-\frac14.
\]
Add $m-f-\nu$ and use Proposition~\ref{prop:decomposition}.
\end{proof}
The strict inequality $m>1/4$ cannot be dropped: a balanced complete bipartite weighted host has $m=1/4$ and $S=E(Q)$, so $\Phi=0$.

For a weighted host write
\[
 \MC(Q)=\max_{U\subseteq V(Q)}c\bigl(E(U,V(Q)\setminus U)\bigr),
 \qquad b(Q)=m-\MC(Q).
\]

\begin{lemma}[Triangle mass from cut deficit]\label{lem:cuttriangle}
For every weighted host,
\begin{equation}\label{eq:cuttriangle}
 \MC(Q)\ge4m^2-2\Delta,
 \qquad 2\Delta\ge b(Q)+4m(m-1/4).
\end{equation}
If $m\ge1/4-\xi$, $b(Q)\ge\rho>0$, and $0\le\xi\le\rho/2$, then $\Delta\ge\rho/4$.
\end{lemma}
\begin{proof}
The cut with one side $N(w)$ has mass
\[
 \sum_{v\in N(w)}x_vD(v)-2\tau(w).
\]
Averaging over $w$, the first term becomes $\sum_vx_vD(v)^2\ge4m^2$, which proves~\eqref{eq:cuttriangle}. If $m\ge1/4$, the last term of~\eqref{eq:cuttriangle} is nonnegative. Otherwise $4m\le1$ and $m-1/4\ge-\xi$, so $4m(m-1/4)\ge-\xi$. In both cases $2\Delta\ge\rho-\xi\ge\rho/2$.
\end{proof}

\begin{theorem}[Stable finite palette bound]\label{thm:stable}
Suppose $\delta>1/3$, $m\ge1/4-\xi$, $b(Q)\ge\rho>0$, and $0\le\xi\le\rho/2$. Then
\begin{equation}\label{eq:stable}
 \Phi\ge\frac32m-\frac14-\frac{2\xi}{\rho}.
\end{equation}
\end{theorem}
\begin{proof}
Lemma~\ref{lem:cuttriangle} gives $\Delta\ge\rho/4$. By the stable assertion of Theorem~\ref{thm:polynomial},
\[
 \frac B\Delta\ge f+\nu+\frac m2-\frac14-\frac{\xi}{2\Delta}
 \ge f+\nu+\frac m2-\frac14-\frac{2\xi}{\rho}.
\]
Now use $\Phi_F\ge B/\Delta$ and~\eqref{eq:decomposition} as before.
\end{proof}
The error term does not depend on the number of vertices of $Q$. This is what allows us to apply the bound after deleting a small number of edges.
